\section{Neurosymbolic Integration}\label{sec:nesy}

This section is where the three pillars meet. The device-resident substrate (\Cref{sec:datalog}) keeps the symbolic state on the GPU; verified knowledge compilation (\Cref{sec:prob}) turns a network's outputs into a sound differentiable circuit; and zero-copy coupling carries gradients back into the neural optimizer without leaving the device. Together they let neural perception and logical reasoning train jointly, end to end, with no host--device transfers in the loop.

\subsection{Neural Predicates}

In xlog a neural network \emph{is} a predicate. The \texttt{nn/4} declaration embeds a network directly into the program:

\begin{lstlisting}[style=xlogstyle,
                   caption={Neural predicate declaration (\texttt{nn/4}).},
                   label={lst:nn-decl}]
nn(network_name, [input_vars], output_var,
   [output_labels]) :: predicate(args).
\end{lstlisting}

This is not a foreign-function escape hatch: it states that evaluating \texttt{predicate(args)} requires a forward pass through the named network, whose softmax over the label set becomes a distribution over probabilistic facts feeding the knowledge-compilation pipeline of \Cref{sec:prob}. On the Python side a PyTorch module is registered against the predicate:

\begin{lstlisting}[style=pythonstyle,
                   caption={Registering a PyTorch module against a neural predicate.},
                   label={lst:nn-register}]
program.register_network("mnist_net", net, optimizer)
\end{lstlisting}

The dataflow runs in two directions. \textbf{Forward:} when evaluation reaches an \texttt{nn/4}-backed predicate, the registered module runs, and its softmax vector is decomposed into weighted probabilistic facts---one per label---fed into the PIR/CNF/D4/XGCF pipeline. Because circuit structure depends on the program, not the weights, the compiled circuit is cached and only the leaf weights are updated on later iterations. \textbf{Backward:} after the circuit computes the forward log-probability and backward gradients via level-parallel kernels, the per-leaf gradient buffers are exported as DLPack tensors and injected into PyTorch's autograd graph, and the optimizer updates the network as usual. The GIL is released during GPU work, so CUDA execution does not block the interpreter.

\subsection{End-to-End Training}

The canonical example is MNIST addition (the DeepProbLog benchmark): given two digit images, predict their sum, with supervision only on sums---the network never sees individual digit labels. The entire reasoning structure is two rules:

\begin{lstlisting}[style=xlogstyle,
                   caption={MNIST addition: a CNN-backed digit predicate and an addition rule.},
                   label={lst:mnist-prog}]
nn(mnist_net, [X], Y, [0,1,2,3,4,5,6,7,8,9])
    :: digit(X, Y).
addition(X, Y, Z) :-
    digit(X, D1), digit(Y, D2), Z is D1 + D2.
\end{lstlisting}

The \texttt{nn/4} declaration connects the CNN to \texttt{digit/2}; the addition rule computes every consistent decomposition of a sum through probabilistic marginalization. Each query \texttt{addition(i, j, s)} trains by minimizing $-\log P(\texttt{addition}(i, j, s))$ under the compiled circuit. The driver compiles the program, registers the network, and runs the loop, which batches queries sharing a circuit template and runs forward/backward over the cached circuit:

\begin{lstlisting}[style=pythonstyle,
                   caption={Python driver: compile, register, train.},
                   label={lst:mnist-train}]
program = pyxlog.Program.compile(SRC)
program.register_network("mnist_net", net, optimizer)
program.add_tensor_source("train", train_images)
history = pyxlog.train_model_tensor(
    program, queries,
    epochs=epochs, batch_size=batch_size)
\end{lstlisting}

\begin{figure*}[tbp]
  \centering
  \resizebox{0.86\linewidth}{!}{%
  \begin{tikzpicture}
    \node[fneural] (net) {neural net\\(PyTorch)};
    \node[fproc,right=30mm of net] (sym) {symbolic\\reasoning};
    \node[fcert,right=14mm of sym] (circ) {verified circuit\\(compiled once)};
    \node[fnode,right=18mm of circ] (loss) {loss\\$-\log P$};
    \draw[farrow] (net) -- node[flabel,above,align=center]{softmax $\to$\\weighted facts} (sym);
    \draw[farrow] (sym) -- (circ);
    \draw[farrow] (circ) -- node[flabel,above]{forward} (loss);
    \draw[fgrad] (loss.south) to[out=-90,in=-90,looseness=0.7]
          node[flabel,below]{gradients via DLPack (zero-copy)} (net.south);
    \begin{scope}[on background layer]
      \node[fill=figband,rounded corners,fit=(sym)(circ),inner sep=3mm] (band) {};
    \end{scope}
    \node[flabel,figblue,anchor=south] at (band.north) {GPU device-resident};
  \end{tikzpicture}}
  \caption{One end-to-end neurosymbolic iteration, exhibiting the three pillars: a device-resident symbolic substrate (the shaded region never leaves the GPU), verified knowledge compilation (the circuit is compiled and proven correct once, then reused), and zero-copy differentiable coupling (gradients return through DLPack into PyTorch's autograd graph).}
  \label{fig:training}
\end{figure*}

The decisive performance property (\Cref{fig:training}) is that compilation and verification happen once, on the first forward pass; every later iteration reuses the cached circuit, executing only the weight update and the level-parallel forward/backward kernels. This is the mechanism behind the $2.74\times$ training speedup; the loss reduction across seeds confirms that gradients flow from circuit evaluation through the logic program back to the CNN, and the full timing breakdown is reported in \Cref{sec:eval}.

\subsection{Zero-Copy Interoperability}\label{subsec:interop}

A GPU-native engine is only useful if results reach the frameworks researchers already use---without copies that would reintroduce the transfer wall. xlog exposes its device-resident state through four standard interfaces.

\textbf{DLPack.} Each relation column or gradient buffer becomes a managed DLPack~\cite{dlpack} tensor whose GPU memory is reference-counted and freed only when every consumer releases it; on import, xlog validates device, dtype, contiguity, and alignment, optionally against an expected schema. In Python the protocol surfaces as a capsule, so \texttt{torch.from\_dlpack} yields a live CUDA tensor backed by the very memory xlog wrote.

\textbf{Arrow.} For columnar tools (Polars, DuckDB, cuDF), relations export through the Arrow~\cite{arrow} C Device Data Interface; symbol columns export as Arrow dictionary arrays, so a consumer reads symbolic columns as native string dictionaries while xlog keeps the compact integer representation.

\textbf{Python bindings.} The \texttt{pyxlog} extension (PyO3~\cite{pyo3}) exposes compilation, evaluation, training, and result extraction; tensor-valued results are returned as DLPack capsules, and long-running GPU operations release the GIL so data-loader threads proceed concurrently.

\textbf{PyTorch autograd.} The circuit behaves as a differentiable function inside PyTorch: the network's grad-enabled forward feeds the circuit, the scalar loss is exported via DLPack and re-imported, and a batched path stacks inputs into one forward pass and runs a single backward---so gradients flow from the logic layer back to \texttt{nn.Module} parameters with no custom autograd \texttt{Function}.

\subsection{Term Embeddings and Differentiable ILP}

Two further paths ride the same device-resident bridge. \textbf{Term embeddings} attach dense, optionally trainable vectors to logical symbols: where \texttt{nn/4} maps perception to discrete facts, embeddings give logical entities a continuous representation that participates in neural computation, with gradients flowing through the embedding lookup---enabling, for example, knowledge-graph entity representations trained jointly with neural perception. \textbf{Differentiable ILP} learns first-order rules rather than network weights: candidate rules are represented as sparse bitmasks over the ground-atom space, credit assignment runs entirely on-device via scatter--gather, and a six-gate promotion pipeline (convergence, novel-rate, regression, held-out F1, ambiguity, typed-schema checks) controls which rules are committed. This sparse, GPU-resident formulation avoids the cubic blowup of dense rule materialization; it is a beta subsystem (\Cref{sec:limits}).
